\documentclass[12pt]{article}
\usepackage[utf8]{inputenc}
\usepackage[T1]{fontenc}
\usepackage[margin=1in]{geometry}
\usepackage{setspace}
\usepackage{xcolor}
\usepackage[colorlinks=true, urlcolor=blue!60!black, linkcolor=blue!60!black]{hyperref}
\setlength{\parindent}{0em}
\setlength{\parskip}{0.45em}
\renewcommand{\baselinestretch}{1.0}
\begin{document}
\begin{center}
{\Large \textit{Romer v. Evans}, 517 U.S. 620 (1996)\par}
\vspace{0.3em}
{\normalsize Case Brief --- One-Page Edition\par}
\end{center}

\textbf{Memory Jogger.} The Colorado Amendment 2 case: animus fails rational-basis review (\href{https://www.courtlistener.com/opinion/118027/romer-v-evans/}{517 U.S. 620}).

\textbf{Facts.} In 1992 Colorado voters approved Amendment 2, repealing Aspen, Boulder, and Denver ordinances and state policies banning sexual-orientation discrimination, and barring every level of government from adopting new ones.

\textbf{Procedural History.} The Colorado courts enjoined it under strict scrutiny; the Supreme Court granted certiorari.

\textbf{Votes.} Affirmed, 6--3. Kennedy wrote for the Court; Scalia, Rehnquist, and Thomas dissented.

\textbf{Holding.} Amendment 2 violates equal protection: a broad disability imposed on a single class, unconnected to any legitimate purpose, reflects a bare desire to harm a politically unpopular group.

\textbf{Reasoning.} Amendment 2 imposed a ``special disability upon those persons alone,'' not a denial of special rights, and was ``at once too narrow and too broad,'' raising ``the inevitable inference'' of animosity. Scalia saw a defense of traditional mores, rational under \textit{Bowers v. Hardwick}, and decried ``political will'' over judgment.

\textbf{Examples.} \emph{Same-Side:} a ban on protections for unmarried cohabitants. \emph{Opposite-Side:} a uniform statewide employment code. \emph{Fence-Sitter:} a single city's charter amendment.

\textbf{Critique.} Critics fault the silence on \textit{Bowers} and an unannounced ``rational basis with bite,'' while defenders recast it as a pariah principle or a bill of attainder.

\textbf{Key Quotation.} ``A law declaring that in general it shall be more difficult for one group of citizens than for all others to seek aid from the government is itself a denial of equal protection of the laws in the most literal sense.'' --- Justice Kennedy, Opinion of the Court, 517 U.S. at 633.

\end{document}
